\subsection{Normal subgroups}
In this section, we will define the concept of normal subgroup and we will 
also talk about how normal subgroups are related with the concept of quotient group, 
and conjugacy classes. We will also define the concept of simple group and we will
discuss the Lagrange's theorem as a mean to help us prove that a given group is simple.

\begin{definition}[Normal subgroup]
    For every $H$ subgroup of $G$, we say that $H$ is a normal subgroup of $G$, and we indicate it 
    with $H \triangleleft G$ if it's closed under conjugation:

    $$ H \triangleleft G \iff \forall h \in H\;\;\forall g \in G\;\; ghg^{-1} \in H $$
    $$ H \triangleleft G \iff \forall g \in G\;\; gHg^{-1} = H $$
\end{definition} \vspace{0.2em}

\begin{lemma} 
\label{lem:normal_cosets}
    If a group $G$ has a normal subgroup $H$, then the left cosets of $H$ in $G$ are equal to 
    the right cosets of $H$ in $G$ and viceversa.
\end{lemma} 
\begin{proof}
    $$ gHg^{-1} = H $$
    $$ gHg^{-1}g = Hg $$
    $$ gH = Hg $$
\end{proof} \vspace{0.2em}

\begin{definition}[Conjugacy class]
    The conjugacy class of an element $g \in G$ is the set of all elements that are conjugate to $g$ 
    in $G$, that is, the set of all $hgh^{-1}$: 
    
    $$ \texttt{Conj}_{G}(g) = \{ hgh^{-1} \mid h \in G \} $$
\end{definition} \vspace{0.2em}

\begin{lemma}
\label{lem:conjugacy_classes}
    Conjugancy classes are either completly disjoint or completly identical
\end{lemma}
\begin{proof}
    
    $$ \texttt{Conj}_{G}(a) \cap \texttt{Conj}_{G}(b) \neq \varnothing \Rightarrow 
        [\; \exists h,k \in G\;\; hah^{-1} = kbk^{-1} \;] $$
    $$ hah^{-1} = kbk^{-1} \Rightarrow k^{-1}hah^{-1} = bk^{-1} \Rightarrow (k^{-1}h)a(h^{-1}k) = b $$
    $$ hah^{-1} = kbk^{-1} \Rightarrow ah^{-1} = h^{-1}kbk^{-1} \Rightarrow a = (h^{-1}k)b(k^{-1}h) $$
\end{proof} \vspace{0.2em}

\begin{theorem}[Lagrange's theorem]
\label{thm:lagrange}
    The order of $H$ subgroup of $G$ divides $|G|$.
\end{theorem}
\begin{proof}
    The left cosets in $H$, meaning the family of sets $gH$ for some fixed $g \in G$, 
    are all of the same size, since the value of $g$ doesn't alter the number 
    of elements in $H$. More specifically, the size of the left cosets is equal 
    to the size of $H$. 

    $$ a \sim b \iff a \in bH $$ \vspace{0.2em}

    We can easily see that the $\sim$ relation is an equivalence
    relation, thus, $gH$ are a family of equivalence classes, herby they partion $G$, hence  
    the order of $G$ is the sum of the orders of the left cosets (which all have the same size):

    $$ |G| = |g_1H| + |g_2H| + \dots = |H| + |H| + \dots = n \cdot |H| $$

\end{proof}

\newpage

\begin{lemma}
\label{lem:intersectionOfNormalIsNormal}
    The intersection of two normal subgroups is a normal subgroup.
\end{lemma}
\begin{proof}
    Let $H_1$ and $H_2$ be two normal subgroups of $G$. Then, by definition of normality:

    $$ \forall h \in (H_1 \cap H_2) \;\;\forall g \in G\;\; ghg^{-1} \in H_1 \land ghg^{-1} \in H_2 $$
    $$ \therefore\;\; ghg^{-1} \in (H_1 \cap H_2) $$
\end{proof} \vspace{0.2em}

\begin{definition}[Simple group]
    A group $G$ is called simple if it has no non-trivial normal subgroups. In other words, the only 
    normal subgroups of $G$ are itself and $\{\epsilon\}$.
\end{definition} \vspace{0.2em}

\begin{definition}[Quotient of groups]
    The quotient $\frac{G}{H}$ is the set of all the left cosets of $H$ in $G$. 
    By this, we mean that for every fixed $g \in G$, $gH$ must be an element of $\frac{G}{H}$.
\end{definition} \vspace{0.2em}

We are intrested into studying the quotient of groups, in particular, we 
want to understand when the quotient of a groups is a group itself. To do so, we will first 
need to define an operation on elements of a quotient. Such operation is $\times$: \\

$$ aH \times bH = (ab)H $$ \vspace{0.2em}

This operation is well defined only if for every possible choice of $n_1, n_2 \in N$ 
exists and $n_3 \in N$ such that the following holds true: \\

$$ (an_1)(bn_2) = (ab)n_3 $$
$$ (an_1)(bn_2) \in (ab)N $$ \vspace{0.2em}

\begin{theorem}
\label{thm:quotientGroup}
    The quotient $(\frac{G}{H}, \times)$ is a group if and only if $H$ is normal subgroup of $G$
\end{theorem}
\begin{proof}
    $$ (an_1)(bn_2) \in (ab)N $$
    $$ (an_1)(bn_2) = (ab)n_3 $$
    $$ an_1bn_2 = abn_3 $$
    $$ n_1bn_2 = bn_3 $$
    $$ b^{-1}n_1bn_2 = n_3 $$
    $$ b^{-1}n_1b = n_3n_2^{-1} $$
    $$ b^{-1}n_1b \in N $$
    $$ b^{-1}Nb = N $$
\end{proof} \vspace{0.2em}

\newpage